\begin{definition}
Let $n\ge2$ and $p>1$.  The predicate
$\mathrm{CitedHolderOscillationMinimum}(n,p)$ is part~(a) of the cited result
(C4) of the Note for the formalization (De
Giorgi--Ladyzhenskaya--Ural'tseva H\"older estimates;
[G, chapter VII, Theorems 2.1--2.2], [GG2], [LU, chapter 4]):

\emph{for all $0<\lambda\le\Lambda$, $a>0$ and
$M\ge0$ there are $\alpha\in(0,1]$ and $d>0$ such that, for every
$\varepsilon\in[0,1]$: if $\Omega$ is open, $f$ satisfies (F1) of the thesis on
$\Omega$ with these parameters (\noderef{VariationalGrowthF1}), $u$ is a local
minimum for $F$ in $\Omega$ (\noderef{IsLocalMinimumOn}) with $|u|\le M$ a.e.\ on
$\Omega$, and $\overline B_R(x_0)\subseteq\Omega$, then for $0<\rho\le R$}
\[
\operatorname*{osc}_{B_\rho}u\;\le\;4^\alpha\max\Bigl\{\operatorname*{osc}_{B_R}u,\;dR\Bigr\}\Bigl(\frac\rho R\Bigr)^\alpha .
\]

This is the estimate displayed on thesis p.~31.  The constants $\alpha$ and $d$
are bound after $n,p$, the structure constants and $M=\|u\|_{L^\infty}$, and
before $\varepsilon$, $\Omega$, the integrand and $u$.  That is exactly the
dependence the Note allows: at l.~263 it says the constants ``may be allowed to
depend also on $M=\|u\|_{L^\infty}$'', naming $M$ and nothing else, so
$\varepsilon$ is not among the quantities $\alpha$ and $d$ may depend on and
one pair of constants must serve the whole range $0\le\varepsilon\le1$.

\medskip
\noindent\textbf{Why this is a predicate of its own and not a field of a
two-part bundle.}  The Note at l.~263 itself labels the two halves of (C4) as
``(a) A bounded local minimum \ldots'' and ``(b) a bounded weak solution
\ldots'', and they share no parameter beyond $n$ and $p$: (a) is a statement
about a variational integrand under (F1) of the thesis, (b) one about an operator
pair under (Q1) of the thesis--(Q5) of the thesis.  Carrying them as two fields of one
structure has a cost that is invisible in the mathematics and decisive in the
formalisation: a consumer of (a) must then assume the conjunction, so if (b)
were false as a $\mathrm{Prop}$ --- which it is if the operator pair is allowed
to be non-Carath\'eodory, the defect repaired in
\noderef{EllipticStructureQ12} and \noderef{EllipticStructureQ345} --- then
every statement assuming the bundle would be vacuously true, including the four
that use only (a).  Splitting the two confines any such failure to the half that
has it.  The part (b) predicate is \noderef{CitedHolderOscillationSolution}.

\medskip
\noindent\textbf{Why $\varepsilon$ is bound after the two constants, and not
before.}  This is not a strengthening beyond the literature; it is what the
literature states and what \S4 needs.  The De Giorgi estimates are proved from
the growth condition (F1) of the thesis (resp.\ (Q1) of the thesis--(Q5) of the thesis),
whose weight is sandwiched uniformly over the range in question:
\[
|h|^{p}\;\le\;\bigl(\varepsilon+|h|^{2}\bigr)^{\frac p2}\;\le\;\bigl(1+|h|^{2}\bigr)^{\frac p2}
\qquad(0\le\varepsilon\le1),
\]
so the constants the iteration produces depend on the two ends of that
sandwich, not on where $\varepsilon$ sits inside it.  The ordering is also
forced by the use: Lemma~Lemma~4.1 of the thesis's own list of what $\sigma_2$ and $R_1$
depend on (thesis l.~1163--1164) is $n,p,\lambda,\Lambda,a,\nu,M$ and
$d(\Omega',\partial\Omega)$, with no $\varepsilon$, while $\sigma_0$ and
$\sigma_1$ carry $\sigma(R+C_1R^{\alpha})$ and $C_1R^{\alpha}$ and so make
$\sigma_2$ inherit $R^{\alpha}$.  With $\varepsilon$ bound before the
existential, $\alpha$ would be a function $\alpha(\varepsilon)$ and a single
$\sigma_2$ would have to dominate $R^{\alpha(\varepsilon)}$ for every
$\varepsilon$; since $R^{\alpha}\to1$ as $\alpha\downarrow0$ for each fixed
$R\in(0,1)$, the pointwise supremum of that family is identically $1$ on
$(0,1)$, which no function with $\sigma_2(0)=0$ and $\sigma_2$ increasing can
dominate.  So the earlier ordering made Lemma~Lemma~4.1 of the thesis unprovable rather
than merely harder.  All
oscillations are essential oscillations (\noderef{EssOscOn}); the hypothesis
$|u|\le M$ a.e.\ makes them well defined.  This result is used as an explicit
hypothesis and is never introduced as an axiom.
\end{definition}
